\section{Running and packaging on Windows and macOS}

The code is pure Python (Tk, NumPy/SciPy/pandas, Astropy, Matplotlib,
astroquery) with no compiled components of its own and no platform
conditionals; porting is an installation exercise.

\subsection{Windows}

\begin{enumerate}
  \item Install CPython 3.10+ from \url{python.org} (the installer bundles
        Tk; tick ``Add python.exe to PATH''). Avoid the Microsoft Store
        build, whose sandbox complicates file dialogs and per-user paths.
  \item \code{py -m pip install -r requirements.txt}
  \item Run with \code{py etc\_gui.py} from the distribution folder.
\end{enumerate}

Notes: paths are handled with \code{pathlib}, so backslashes are transparent;
configuration files are written next to the program --- install under a
user-writable directory, not \code{Program Files}; on high-DPI displays Tk
scales automatically on Python $\geq$ 3.12, otherwise call
\code{ctypes.windll.shcore.SetProcessDpiAwareness(1)} early if the UI looks
blurry. For a self-contained distribution use PyInstaller:

\begin{verbatim}
pip install pyinstaller
pyinstaller --windowed --name SPETC ^
    --add-data "data;data" ^
    --add-data "etc_sites.json;." ^
    --add-data "observatories.json;." etc_gui.py
\end{verbatim}

Matplotlib and Astropy need no hooks beyond PyInstaller's built-ins; test
the frozen build offline, because Astropy must fall back to its bundled
IERS table (the code already sets \code{auto\_download = False}).

\subsection{macOS}

\begin{enumerate}
  \item Install CPython from \url{python.org} (universal2; bundles a
        private Tcl/Tk 8.6 --- do \emph{not} rely on the Apple system
        Python). Homebrew works too: \code{brew install python-tk}.
  \item \code{python3 -m pip install -r requirements.txt}
  \item \code{python3 etc\_gui.py}
\end{enumerate}

Notes: on Retina displays Tk renders sharp with the python.org build;
Gatekeeper will quarantine a PyInstaller app (\code{--windowed} produces
\code{SPETC.app}) unless it is signed and notarized --- for personal use,
right-click~$\to$~Open once, or clear the quarantine attribute with
\code{xattr -dr com.apple.quarantine SPETC.app}. The case-sensitive file
systems sometimes used on macOS are already handled: all data files are
referenced with their exact case.

\subsection{Both platforms}

Keep \code{data/} beside the executable (the shipped data-directory entry
is resolved relative to \code{etc\_gui.py}); time zones use the standard
\code{zoneinfo} database (on Windows, \code{pip install tzdata} is pulled
in automatically as a dependency of the requirements); and verify a port
with the two test programs before first use --- they exercise the whole
physics stack without opening a window.
